% !TEX root = ../main.tex
\corrigesChapitre{16}{Énergie potentielle et énergie mécanique}
\corrigesSouples

\begin{blocLivre}{cor-16-1}
\begin{corrige}{1}{Calculer des énergies potentielles}
\demarche{$E_\mathrm{p}=m\,g\,z$, $z$ étant l'altitude par rapport à l'origine choisie (axe vers le haut).}
\begin{questions}
\item $E_\mathrm{p}=0{,}80\times9{,}81\times2{,}0\approx\res{\qty{15,7}{J}}$.
\item Le sol est alors à $z=\qty{-2,0}{m}$ : $E_\mathrm{p}\approx\res{\qty{-15,7}{J}}$. Une énergie potentielle peut être négative.
\item $E_\mathrm{p}=70\times9{,}81\times4808\approx\res{\qty{3,30e6}{J}}$.
\item $\Delta E_\mathrm{p}=70\times9{,}81\times(4808-1035)\approx\res{\qty{2,59e6}{J}}$ ; une variation \res{\text{ne dépend pas}} de l'origine choisie.
\end{questions}
\retiens{Seules les variations d'énergie potentielle ont un sens physique : l'origine est un choix, à garder tout au long de l'exercice.}
\end{corrige}
\end{blocLivre}
\begin{blocLivre}{cor-16-2}
\begin{corrige}{2}{Une chute sans frottement}
\demarche{Seul le poids travaille : $E_\mathrm{m}$ se conserve. $E_\mathrm{c}=E_\mathrm{m}-E_\mathrm{p}$ en chaque point.}
\begin{questions}
\item $E_\mathrm{m}=0+0{,}20\times9{,}81\times10\approx\res{\qty{19,6}{J}}$.
\item La seule force, le poids, est \res{\text{conservative}} : pas de travail de forces non conservatives.
\item $E_\mathrm{p}=0{,}20\times9{,}81\times4{,}0\approx\res{\qty{7,85}{J}}$ ; $E_\mathrm{c}\approx19{,}6-7{,}85\approx\res{\qty{11,8}{J}}$ ; $v=\sqrt{\dfrac{2\times11{,}8}{0{,}20}}\approx\res{\qty{10,8}{m.s^{-1}}}$.
\item Au sol, $E_\mathrm{c}=\qty{19,6}{J}$ : $v=\sqrt{\dfrac{2\times19{,}6}{0{,}20}}\approx\res{\qty{14,0}{m.s^{-1}}}$.
\end{questions}
\end{corrige}
\end{blocLivre}
\begin{blocLivre}{cor-16-3}
\begin{corrige}{3}{Lire des courbes d'énergie}
\demarche{$E_\mathrm{p}$ suit l'altitude (maximale au sommet), $E_\mathrm{c}$ s'annule au sommet, $E_\mathrm{m}$ est leur somme.}
\begin{questions}
\item 1 : $\res{E_\mathrm{p}}$ (maximale au sommet) ; 2 : $\res{E_\mathrm{c}}$ (nulle au sommet) ; 3 : $\res{E_\mathrm{m}}$, somme des deux.
\item $E_\mathrm{p}=0{,}10\times9{,}81\times1{,}0\approx\res{\qty{0,981}{J}}$ ; $E_\mathrm{c}=\dfrac{1}{2}\times0{,}10\times8{,}0^2=\res{\qty{3,2}{J}}$ ; $E_\mathrm{m}\approx\res{\qty{4,18}{J}}$.
\item $E_\mathrm{c}=0$ vers \res{\qty{0,82}{s}} ; alors $m\,g\,z=E_\mathrm{m}$ : $z=\dfrac{4{,}18}{0{,}10\times9{,}81}\approx\res{\qty{4,26}{m}}$.
\item \res{\text{Oui}} : l'énergie mécanique reste constante.
\end{questions}
\schema[corrige]{16-ex03}
\end{corrige}
\end{blocLivre}
\begin{blocLivre}{cor-16-4}
\begin{corrige}{4}{Le pendule simple}
\demarche{Altitude perdue par géométrie, puis conservation de l'énergie mécanique : $\dfrac{1}{2}\,m\,v_\mathrm{B}^2=m\,g\,h$.}
\begin{questions}
\item En A, la bille est à $L\cos\ang{30}$ sous le point d'attache, en B à $L$ : $h=L\,(1-\cos\ang{30})\approx\res{\qty{0,134}{m}}$ (schéma).
\item Le poids et la tension du fil ; la tension, perpendiculaire à la trajectoire, \res{\text{ne travaille pas}}.
\item $v_\mathrm{B}=\sqrt{2\,g\,h}=\sqrt{2\times9{,}81\times0{,}134}\approx\res{\qty{1,62}{m.s^{-1}}}$.
\item $h=1{,}0\times(1-\cos\ang{60})=\qty{0,50}{m}$ : $v_\mathrm{B}=\sqrt{2\times9{,}81\times0{,}50}\approx\res{\qty{3,13}{m.s^{-1}}}$.
\end{questions}
\schema[corrige]{16-ex04}
\end{corrige}
\end{blocLivre}
\begin{blocLivre}{cor-16-5}
\begin{corrige}{5}{Conservation ou non ?}
\demarche{$\Delta E_\mathrm{m}=W(\text{forces non conservatives})$ : il faut repérer ces forces et le signe de leur travail.}
\begin{questions}
\item \res{\text{Conservation}} : seul le poids agit.
\item \res{\text{Diminution}} : $E_\mathrm{c}$ constante, $E_\mathrm{p}$ diminue ; les frottements de l'air dissipent l'énergie.
\item \res{\text{Augmentation}} : $E_\mathrm{c}$ constante, $E_\mathrm{p}$ augmente ; le travail du téléski est moteur.
\item \res{\text{Diminution}} : les frottements dissipent peu à peu l'énergie mécanique.
\end{questions}
\end{corrige}
\end{blocLivre}
\begin{blocLivre}{cor-16-6}
\begin{corrige}{6}{La balle de ping-pong, en Python}
\demarche{$E_\mathrm{m}=\dfrac{1}{2}\,m\,v^2+m\,g\,z$ calculée à chaque date ; sa diminution mesure le travail des frottements.}
\begin{questions}
\item L'\res{\text{énergie mécanique}}, en millijoules ; l'origine des altitudes est le \res{\text{sol}} ($z=0$ en bas).
\item Elle \res{\text{diminue}}, de plus en plus vite : les frottements de l'air, qui croissent avec la vitesse, dissipent de l'énergie.
\item $W(\vec{f})=\Delta E_\mathrm{m}=45{,}3-53{,}0=\res{\qty{-7,7}{mJ}}$, environ \qty{15}{\percent} de l'énergie initiale.
\item $g\,t=9{,}81\times0{,}6\approx\qty{5,89}{m.s^{-1}}$ contre \qty{5,17}{m.s^{-1}} mesurée : la balle, très légère, est nettement \res{\text{freinée par l'air}}.
\end{questions}
\end{corrige}
\end{blocLivre}
\begin{blocLivre}{cor-16-7}
\begin{corrige}{7}{Les montagnes russes}
\demarche{Sans frottement, $E_\mathrm{m}$ se conserve : $v^2=v_\mathrm{A}^2+2\,g\,(z_\mathrm{A}-z)$. Avec frottement, l'énergie dissipée vaut $E_\mathrm{m}(\mathrm{A})-E_\mathrm{m}(\mathrm{B})$.}
\begin{questions}
\item $E_\mathrm{m}=\dfrac{1}{2}\times500\times2{,}0^2+500\times9{,}81\times40\approx\res{\qty{1,97e5}{J}}$.
\item $v_\mathrm{B}=\sqrt{4+2\times9{,}81\times35}\approx\res{\qty{26,3}{m.s^{-1}}}$ ; $v_\mathrm{C}=\sqrt{4+2\times9{,}81\times15}\approx\res{\qty{17,3}{m.s^{-1}}}$.
\item $E_\mathrm{dissipée}=\dfrac{1}{2}\times500\times(26{,}3^2-24^2)\approx\res{\qty{2,87e4}{J}}$.
\item $\dfrac{\num{2,87e4}}{\num{1,97e5}}\approx\res{\qty{15}{\percent}}$ de l'énergie mécanique initiale.
\end{questions}
\schema[corrige]{16-ex07}
\end{corrige}
\end{blocLivre}
\begin{blocLivre}{cor-16-8}
\begin{corrige}{8}{Les rebonds d'une balle}
\demarche{Aux sommets, la vitesse est nulle : $E_\mathrm{m}=m\,g\,h$. Comparer les énergies revient à comparer les hauteurs.}
\begin{questions}
\item Au départ, $E_\mathrm{m}=m\,g\times1{,}20$ ; au sommet du rebond, $m\,g\times0{,}80$ : elle a \res{\text{diminué}}.
\item $\Delta E=\num{0,058}\times9{,}81\times0{,}40\approx\res{\qty{0,228}{J}}$, dissipés en énergie \res{\text{thermique}} (déformation de la balle) et en son.
\item $\dfrac{0{,}80}{1{,}20}\approx\res{\num{0,67}}$, soit deux tiers.
\item $1{,}20\times\left(\dfrac{2}{3}\right)^5\approx\res{\qty{0,158}{m}}$.
\end{questions}
\end{corrige}
\end{blocLivre}
\begin{blocLivre}{cor-16-9}
\begin{corrige}{9}{Une balle lancée vers le haut}
\demarche{Origine au point de lancer. Sans frottement : $\dfrac{1}{2}\,m\,v_0^2=m\,g\,h$. Avec frottement : $W(\vec{f})=\Delta E_\mathrm{m}$.}
\begin{questions}
\item $h=\dfrac{v_0^2}{2\,g}=\dfrac{144}{19{,}62}\approx\res{\qty{7,34}{m}}$.
\item $W(\vec{f})=m\,g\,h_\mathrm{réel}-\dfrac{1}{2}\,m\,v_0^2=0{,}15\times9{,}81\times(6{,}5-7{,}34)\approx\res{\qty{-1,24}{J}}$.
\item $E_\mathrm{m}$ se conserve : la balle repasse à \res{\qty{12}{m.s^{-1}}}.
\item \res{\text{Plus petite}} : les frottements dissipent aussi de l'énergie à la descente.
\end{questions}
\end{corrige}
\end{blocLivre}
\begin{blocLivre}{cor-16-10}
\begin{corrige}{10}{Le téléski}
\demarche{Vitesse constante : $\Delta E_\mathrm{c}=0$, donc $\Delta E_\mathrm{m}=\Delta E_\mathrm{p}=m\,g\,\Delta z$ ; c'est le travail du téléski (sans frottement).}
\begin{questions}
\item $E_\mathrm{c}$ \res{\text{constante}} ; $E_\mathrm{p}$ et $E_\mathrm{m}$ \res{\text{augmentent}}.
\item $\Delta E_\mathrm{m}=75\times9{,}81\times300\approx\res{\qty{2,21e5}{J}}$.
\item $W=\res{\qty{2,21e5}{J}}$ ; $P=\dfrac{\num{2,21e5}}{360}\approx\res{\qty{613}{W}}$.
\item Il doit aussi compenser le travail \res{\text{résistant}} des frottements des skis sur la neige.
\end{questions}
\end{corrige}
\end{blocLivre}
\begin{blocLivre}{cor-16-11}
\begin{corrige}{11}{Le pendule réel}
\demarche{Origine des altitudes au point bas : $E_\mathrm{m}=m\,g\,L\,(1-\cos\theta_\mathrm{max})$ aux positions extrêmes, où $E_\mathrm{c}=0$.}
\begin{questions}
\item 1 : $E_\mathrm{p}$ ; 2 : $E_\mathrm{c}$ ; 3 : $E_\mathrm{m}$. À chaque période, le pendule passe \res{\text{deux fois}} par le bas et deux fois par une position extrême.
\item $E_\mathrm{m}=0{,}20\times9{,}81\times1{,}0\times(1-\cos\ang{40})\approx\res{\qty{0,459}{J}}$.
\item $E_\mathrm{m}=0{,}20\times9{,}81\times(1-\cos\ang{30})\approx\res{\qty{0,263}{J}}$ ; énergie dissipée : $\res{\qty{0,196}{J}}$.
\item $P=\dfrac{0{,}196}{20}\approx\res{\qty{9,8e-3}{W}}$, soit environ \qty{10}{mW}.
\end{questions}
\schema[corrige]{16-ex11}
\end{corrige}
\end{blocLivre}
\begin{blocLivre}{cor-16-12}
\begin{corrige}{12}{La centrale hydroélectrique}
\demarche{Chaque seconde, \qty{50}{m^3} d'eau perdent leur énergie potentielle ; la puissance électrique en est \qty{90}{\percent}.}
\begin{questions}
\item $m=\qty{1000}{kg}$ : $\Delta E_\mathrm{p}=1000\times9{,}81\times300\approx\res{\qty{2,94e6}{J}}$.
\item $P=50\times\num{2,94e6}\approx\res{\qty{1,47e8}{W}}$ ; $P_\mathrm{e}=0{,}90\times\num{1,47e8}\approx\res{\qty{1,32e8}{W}}$ (\qty{132}{MW}).
\item $E=132\times24\approx\res{\qty{3,18e3}{MWh}}$.
\item Par an : $3{,}18\times10^{3}\times365\approx\qty{1,16e6}{MWh}$, soit $\dfrac{\num{1,16e9}}{4700}\approx\res{\num{250000}}$ foyers.
\end{questions}
\end{corrige}
\end{blocLivre}
\begin{blocLivre}{cor-16-13}
\begin{corrige}{13}{Le lancer du poids}
\demarche{Sans frottement, $E_\mathrm{m}$ se conserve entre le lâcher, le sommet et l'arrivée.}
\begin{questions}
\item $E_\mathrm{c}=\dfrac{1}{2}\times7{,}26\times13^2\approx\qty{613}{J}$ ; $E_\mathrm{p}=7{,}26\times9{,}81\times2{,}1\approx\qty{150}{J}$ ; $E_\mathrm{m}\approx\res{\qty{763}{J}}$.
\item Au sol, $E_\mathrm{c}=\qty{763}{J}$ : $v=\sqrt{\dfrac{2\times763}{7{,}26}}\approx\res{\qty{14,5}{m.s^{-1}}}$.
\item $v_\mathrm{sommet}=13\times\cos\ang{40}\approx\qty{9,96}{m.s^{-1}}$ : $E_\mathrm{c}=\dfrac{1}{2}\times7{,}26\times9{,}96^2\approx\res{\qty{360}{J}}$.
\item $E_\mathrm{p}=763-360=\qty{403}{J}$, d'où $z=\dfrac{403}{7{,}26\times9{,}81}\approx\res{\qty{5,66}{m}}$.
\end{questions}
\schema[corrige]{16-ex13-cor}
\end{corrige}
\end{blocLivre}
\begin{blocLivre}{cor-16-14}
\begin{corrige}{14}{La descente à vélo}
\demarche{Vitesse constante : l'énergie potentielle perdue chaque seconde est entièrement dissipée par les frottements.}
\begin{questions}
\item $E_\mathrm{c}$ \res{\text{constante}}, $E_\mathrm{p}$ et donc $E_\mathrm{m}$ \res{\text{diminuent}}.
\item En \qty{1}{s}, elle parcourt \qty{15}{m} sur la route : $\Delta z=15\times0{,}0995\approx\res{\qty{1,49}{m}}$.
\item $P=\dfrac{m\,g\,\Delta z}{\Delta t}=80\times9{,}81\times1{,}49\approx\res{\qty{1,17e3}{W}}$.
\item Descente de $2000\times0{,}0995=\qty{199}{m}$ : $v=\sqrt{15^2+2\times9{,}81\times199}\approx\res{\qty{64,3}{m.s^{-1}}}$, soit \qty{230}{km/h} : les frottements sont indispensables !
\end{questions}
\end{corrige}
\end{blocLivre}
\begin{blocLivre}{cor-16-15}
\begin{corrige}{15}{Le bilan d'énergie des rebonds, en Python}
\demarche{Aux sommets, $E_\mathrm{c}=0$ et $E_\mathrm{m}=m\,g\,h$ : l'énergie mécanique est proportionnelle à la hauteur atteinte.}
\begin{questions}
\item $\dfrac{E_\mathrm{m}(n+1)}{E_\mathrm{m}(n)}=\dfrac{m\,g\,h_{n+1}}{m\,g\,h_n}=\dfrac{h_{n+1}}{h_n}$ : \res{\text{même rapport}}.
\item Elle conserve \qty{73}{\percent} et perd donc environ \res{\qty{27}{\percent}} de son énergie à chaque rebond.
\item Le nombre de rebonds jusqu'à ce que la hauteur passe sous \qty{5}{cm} : \res{\text{11 rebonds}}, la balle semble alors immobile.
\item $\Delta E=\num{0,058}\times9{,}81\times(1{,}50-1{,}10)\approx\res{\qty{0,228}{J}}$.
\end{questions}
\end{corrige}
\end{blocLivre}
